\documentclass[12pt,aspectratio=169,ignorenonframetext]{ltxbeamer}

\usepackage[spanish,es-nodecimaldot,es-tabla]{babel}
\usepackage{multicol}

\title{Bibliografía}
\subtitle{Curso de {\lmr\LaTeX}}
\author{Joar Buitrago\texorpdfstring{\thanks{jebuitragoc@unal.edu.co}}{ <jebuitragoc@unal.edu.co>}}
\date{}

\addbibresource{bibliography.bib}

\newcommand{\BibTeX}{\textsc{Bib}\TeX}



\begin{document}
\begin{frame}
\maketitle
\end{frame}





\section{El problema}

\begin{frame}
\frametitle{¿Por qué bibliografía?}

En cualquier documento escrito muchas veces es necesario incluir bibliografía o referencias.

\pause

\begin{itemize}[<+->]
    \item Se tienen que escribir a mano.
    \item Se debe mantener un estilo o reglas consistentes.
    \item Las citaciones deben mantener un estilo consistente.
    \item Si se cambia el estilo, hay que hacer todo de nuevo.
\end{itemize}
\end{frame}





\begin{frame}
\frametitle{Ejemplo}

\begin{enumerate}
    \lmr
    \setbeamercolor{enumerate item}{fg=white}
    \setbeamertemplate{enumerate item}{[\insertenumlabel]}
    \item García Márquez, Gabriel. Cien años de soledad, 1967.
    \item 1967.  G. G. Márquez. Editorial Sudamericana.
    \item Cien Años de Soledad. (1967) por G. G. Márquez. Editorial Sudamericana.
\end{enumerate}
\end{frame}




\begin{frame}[fragile]
\frametitle{Entorno \texttt{thebibliography}}

    \begin{minted}[breaklines]{latex}
\begin{thebibliography}{9}
    \bibitem{GM1} García Márquez, Gabriel. Cien años de soledad, 1967.
    \bibitem{GM2} 1967.  G. G. Márquez. Editorial Sudamericana.
    \bibitem{GM3} Cien Años de Soledad. (1967) por G. G. Márquez. Editorial Sudamericana.
\end{thebibliography}
    \end{minted}
\end{frame}





\section{{\lmr\BibTeX}}

\begin{frame}
\centering
\vfill
\resizebox{0.95\linewidth}{!}{\lmr\BibTeX}
\vfill
\end{frame}






\begin{frame}
\frametitle{¿Qué es {\lmr\BibTeX}?}

{\lmr\BibTeX} es un sistema automatizado para la gestión bibliográfica.

\pause

Fue creado por Oren Patashnik y Leslie Lamport en 1985.
\end{frame}





\begin{frame}
\frametitle{La base de datos}

Para poder trabajar con {\lmr\BibTeX} necesitamos crear un archivo de base de datos.

\pause

Estos archivos suelen tener la extensión \texttt{.bib}.
\end{frame}






\begin{frame}[fragile]
\frametitle{La base de datos}

\begin{columns}
    \begin{column}{0.55\textwidth}
        \tikz[remember picture,baseline]{\node [anchor=base,inner sep=0] (top) {\strut};}
        \begin{minted}{bibtex}
@book{TeXbook,
    author = {Donald Ervin Knuth},
    title = {The {\lmr\TeX} book},
    publisher = {Stanford University},
    year = {1986},
    edition = {1},
}
        \end{minted}
        \tikz[remember picture,baseline]{\node [anchor=base,inner sep=0] (bottom) {\strut};}
    \end{column}

    \begin{column}{0.45\textwidth}
        \tikz[remember picture,baseline]{\coordinate [anchor=base] (braces);}
    \end{column}
\end{columns}

\begin{tikzpicture}[remember picture, overlay]
    \pause
    \coordinate (keyline) at ($(top)!1/7!(bottom)$);
    \draw [<-] (keyline -| braces) -- ++(1em,0) node [right] {Tipo y Clave};

    \pause
    \node (parameters begin) at ($(top)!2/7!(bottom)$) {\strut};
    \node (parameters end) at ($(top)!5/7!(bottom)$) {\strut};
    \draw [thick,decoration=brace,decorate] (parameters begin.north -| braces) -- (parameters end.south -| braces) node [midway, right] {Parametros};
\end{tikzpicture}
\end{frame}





\begin{frame}[fragile]
\frametitle{La base de datos}

\begin{columns}
    \begin{column}{0.5\textwidth}
        \begin{minted}[escapeinside=||]{latex}
\cite{TeXbook}|\pause|

\nocite{TeXbook}
\nocite{*}|\pause|

\bibliographystyle{plain}
\bibliography{|\textit{documento}|}
        \end{minted}
    \end{column}
    \pause
    {\color{gray} \vline}
    \begin{column}{0.5\linewidth}
        \lmr
        \begin{enumerate}
            \setbeamerfont{enumerate item}{family=\lmr}
            \setbeamercolor{enumerate item}{fg=white}
            \setbeamertemplate{enumerate item}{[\insertenumlabel]}
            \item Donald Knuth. \textit{The {\TeX} book}. Stanford University, 1 edition, 1986.
        \end{enumerate}
    \end{column}
\end{columns}
\end{frame}





\subsection{Compilación}

\begin{frame}
\frametitle{Compilación}

\begin{center}
\begin{tikzpicture}
    \begin{scope}[local bounding box=tex]
        \useasboundingbox (-2ex,-2.5ex) rectangle (2ex,2.5ex);
        \draw [very thick]
            [rounded corners=2pt] (-1.5ex,-2ex) coordinate (element) -- ++(0ex,4ex)
            [rounded corners=1pt] -- ++(2ex,0ex) -- ++(1ex,-1ex)
            [rounded corners=2pt] -- ++(0ex,-3ex) -- cycle
            (element) ++(2ex,4ex) |- ++(1ex,-1ex);
        \draw (element) ++(3ex,0ex) node [fill=black, inner sep=0.5ex] {\texttt{TEX}};
    \end{scope}

    \begin{scope}[local bounding box=compiler, shift={([xshift=2cm]tex)}]
        \alt<6->{\node [black!75] {\lmr\LaTeX};}{\node {\lmr\LaTeX};}
    \end{scope}

    \alt<6->{\draw [->, very thick, black!75] (tex) -- (compiler);}{\draw [->, very thick] (tex) -- (compiler);}

    \begin{scope}[local bounding box=pdf, shift={([xshift=2cm]compiler)}]
        \useasboundingbox (-2ex,-2.5ex) rectangle (2ex,2.5ex);
        \alt<6->{
            \draw [very thick, black!75]
                [rounded corners=2pt] (-1.5ex,-2ex) coordinate (element) -- ++(0ex,4ex)
                [rounded corners=1pt] -- ++(2ex,0ex) -- ++(1ex,-1ex)
                [rounded corners=2pt] -- ++(0ex,-3ex) -- cycle
                (element) ++(2ex,4ex) |- ++(1ex,-1ex);
        }{
            \draw [very thick]
                [rounded corners=2pt] (-1.5ex,-2ex) coordinate (element) -- ++(0ex,4ex)
                [rounded corners=1pt] -- ++(2ex,0ex) -- ++(1ex,-1ex)
                [rounded corners=2pt] -- ++(0ex,-3ex) -- cycle
                (element) ++(2ex,4ex) |- ++(1ex,-1ex);
        }
        \alt<6->{
            \draw [black!75] (element) ++(3ex,0ex) node [fill=black, inner sep=0.5ex] {\texttt{PDF}};
        }{
            \draw (element) ++(3ex,0ex) node [fill=black, inner sep=0.5ex] {\texttt{PDF}};
        }
    \end{scope}

    \alt<6->{\draw [->, very thick, black!75] (compiler) -- (pdf);}{\draw [->, very thick] (compiler) -- (pdf);}

    \begin{scope}[local bounding box=log, shift={([yshift=2cm]pdf)}]
        \useasboundingbox (-2ex,-2.5ex) rectangle (2ex,2.5ex);
        \alt<6->{
            \draw [very thick, black!75]
                [rounded corners=2pt] (-1.5ex,-2ex) coordinate (element) -- ++(0ex,4ex)
                [rounded corners=1pt] -- ++(2ex,0ex) -- ++(1ex,-1ex)
                [rounded corners=2pt] -- ++(0ex,-3ex) -- cycle
                (element) ++(2ex,4ex) |- ++(1ex,-1ex);
        }{
            \draw [very thick]
                [rounded corners=2pt] (-1.5ex,-2ex) coordinate (element) -- ++(0ex,4ex)
                [rounded corners=1pt] -- ++(2ex,0ex) -- ++(1ex,-1ex)
                [rounded corners=2pt] -- ++(0ex,-3ex) -- cycle
                (element) ++(2ex,4ex) |- ++(1ex,-1ex);
        }
        \alt<6->{
            \draw [black!75] (element) ++(3ex,0ex) node [fill=black, inner sep=0.5ex] {\texttt{LOG}};
        }{
            \draw (element) ++(3ex,0ex) node [fill=black, inner sep=0.5ex] {\texttt{LOG}};
        }
    \end{scope}

    \begin{scope}[local bounding box=aux, shift={([yshift=-2cm]pdf)}]
        \useasboundingbox (-2ex,-2.5ex) rectangle (2ex,2.5ex);
        \alt<6->{
            \draw [very thick, black!75]
                [rounded corners=2pt] (-1.5ex,-2ex) coordinate (element) -- ++(0ex,4ex)
                [rounded corners=1pt] -- ++(2ex,0ex) -- ++(1ex,-1ex)
                [rounded corners=2pt] -- ++(0ex,-3ex) -- cycle
                (element) ++(2ex,4ex) |- ++(1ex,-1ex);
        }{
            \draw [very thick]
                [rounded corners=2pt] (-1.5ex,-2ex) coordinate (element) -- ++(0ex,4ex)
                [rounded corners=1pt] -- ++(2ex,0ex) -- ++(1ex,-1ex)
                [rounded corners=2pt] -- ++(0ex,-3ex) -- cycle
                (element) ++(2ex,4ex) |- ++(1ex,-1ex);
        }
        \alt<6->{
            \draw [black!75] (element) ++(3ex,0ex) node [fill=black, inner sep=0.5ex] {\texttt{AUX}};
        }{
            \draw (element) ++(3ex,0ex) node [fill=black, inner sep=0.5ex] {\texttt{AUX}};
        }
    \end{scope}

    \alt<6->{\draw [->, very thick, black!75] (compiler) -- (log);}{\draw [->, very thick] (compiler) -- (log);}
    \alt<6->{\draw [->, very thick, black!75] (compiler) -- (aux);}{\draw [->, very thick] (compiler) -- (aux);}

    \pause

    \begin{scope}[local bounding box=bib,shift={([yshift=2cm]tex)}]
        \useasboundingbox (-2ex,-2.5ex) rectangle (2ex,2.5ex);
        \draw [very thick]
            [rounded corners=2pt] (-1.5ex,-2ex) coordinate (element) -- ++(0ex,4ex)
            [rounded corners=1pt] -- ++(2ex,0ex) -- ++(1ex,-1ex)
            [rounded corners=2pt] -- ++(0ex,-3ex) -- cycle
            (element) ++(2ex,4ex) |- ++(1ex,-1ex);
        \draw (element) ++(3ex,0ex) node [fill=black, inner sep=0.5ex] {\texttt{BIB}};
    \end{scope}

    \alt<6->{\draw [->, very thick, shorten <=0.65em, black!75] (bib) -- (compiler);}{\draw [->, very thick, shorten <=0.65em] (bib) -- (compiler);}

    \pause

    \begin{scope}[local bounding box=bibtex, shift={([xshift=2cm]pdf)}]
        \alt<6->{\node [black!75] {\lmr\BibTeX};}{\node {\lmr\BibTeX};}
    \end{scope}

    \alt<6->{\draw [->, very thick, black!75] (aux) -- (bibtex);}{\draw [->, very thick] (aux) -- (bibtex);}

    \pause

    \begin{scope}[local bounding box=blg, shift={([xshift=2cm]bibtex)}]
        \useasboundingbox (-2ex,-2.5ex) rectangle (2ex,2.5ex);
        \alt<6->{
            \draw [very thick, black!75]
                [rounded corners=2pt] (-1.5ex,-2ex) coordinate (element) -- ++(0ex,4ex)
                [rounded corners=1pt] -- ++(2ex,0ex) -- ++(1ex,-1ex)
                [rounded corners=2pt] -- ++(0ex,-3ex) -- cycle
                (element) ++(2ex,4ex) |- ++(1ex,-1ex);
        }{
            \draw [very thick]
                [rounded corners=2pt] (-1.5ex,-2ex) coordinate (element) -- ++(0ex,4ex)
                [rounded corners=1pt] -- ++(2ex,0ex) -- ++(1ex,-1ex)
                [rounded corners=2pt] -- ++(0ex,-3ex) -- cycle
                (element) ++(2ex,4ex) |- ++(1ex,-1ex);
        }
        \alt<6->{
            \draw [black!75] (element) ++(3ex,0ex) node [fill=black, inner sep=0.5ex] {\texttt{BLG}};
        }{
            \draw (element) ++(3ex,0ex) node [fill=black, inner sep=0.5ex] {\texttt{BLG}};
        }
    \end{scope}

    \begin{scope}[local bounding box=bbl, shift={([yshift=-2cm]blg)}]
        \useasboundingbox (-2ex,-2.5ex) rectangle (2ex,2.5ex);
        \alt<6-7>{
            \draw [very thick, black!75]
                [rounded corners=2pt] (-1.5ex,-2ex) coordinate (element) -- ++(0ex,4ex)
                [rounded corners=1pt] -- ++(2ex,0ex) -- ++(1ex,-1ex)
                [rounded corners=2pt] -- ++(0ex,-3ex) -- cycle
                (element) ++(2ex,4ex) |- ++(1ex,-1ex);
        }{
            \draw [very thick]
                [rounded corners=2pt] (-1.5ex,-2ex) coordinate (element) -- ++(0ex,4ex)
                [rounded corners=1pt] -- ++(2ex,0ex) -- ++(1ex,-1ex)
                [rounded corners=2pt] -- ++(0ex,-3ex) -- cycle
                (element) ++(2ex,4ex) |- ++(1ex,-1ex);
        }
        \alt<6-7>{
            \draw [black!75] (element) ++(3ex,0ex) node [fill=black, inner sep=0.5ex] {\texttt{BBL}};
        }{
            \draw (element) ++(3ex,0ex) node [fill=black, inner sep=0.5ex] {\texttt{BBL}};
        }
    \end{scope}

    \alt<6->{\draw [->, very thick, black!75] (bibtex) -- (bbl);}{\draw [->, very thick] (bibtex) -- (bbl);}
    \alt<6->{\draw [->, very thick, black!75] (bibtex) -- (blg);}{\draw [->, very thick] (bibtex) -- (blg);}

    \pause

    \begin{scope}[local bounding box=compiler2, shift={([xshift=2cm]blg)}]
        \node {\lmr\LaTeX};
    \end{scope}

    \pause
    \pause

    \draw [->, very thick] (tex) -- (compiler2);
    \draw [->, very thick, rounded corners=5pt] (bib) -| (compiler2);
    \pause
    \draw [->, very thick] (bbl) -- (compiler2);

    \pause

    \begin{scope}[local bounding box=pdf2, shift={([xshift=2cm]compiler2)}]
        \useasboundingbox (-2ex,-2.5ex) rectangle (2ex,2.5ex);
        \draw [very thick]
            [rounded corners=2pt] (-1.5ex,-2ex) coordinate (element) -- ++(0ex,4ex)
            [rounded corners=1pt] -- ++(2ex,0ex) -- ++(1ex,-1ex)
            [rounded corners=2pt] -- ++(0ex,-3ex) -- cycle
            (element) ++(2ex,4ex) |- ++(1ex,-1ex);
        \draw (element) ++(3ex,0ex) node [fill=black, inner sep=0.5ex] {\texttt{PDF}};
    \end{scope}

    \draw [->, very thick] (compiler2) -- (pdf2);

    \begin{scope}[local bounding box=log2, shift={([yshift=2cm]pdf2)}]
        \useasboundingbox (-2ex,-2.5ex) rectangle (2ex,2.5ex);
        \draw [very thick]
            [rounded corners=2pt] (-1.5ex,-2ex) coordinate (element) -- ++(0ex,4ex)
            [rounded corners=1pt] -- ++(2ex,0ex) -- ++(1ex,-1ex)
            [rounded corners=2pt] -- ++(0ex,-3ex) -- cycle
            (element) ++(2ex,4ex) |- ++(1ex,-1ex);
        \draw (element) ++(3ex,0ex) node [fill=black, inner sep=0.5ex] {\texttt{LOG}};
    \end{scope}

    \begin{scope}[local bounding box=aux2, shift={([yshift=-2cm]pdf2)}]
        \useasboundingbox (-2ex,-2.5ex) rectangle (2ex,2.5ex);
        \draw [very thick]
            [rounded corners=2pt] (-1.5ex,-2ex) coordinate (element) -- ++(0ex,4ex)
            [rounded corners=1pt] -- ++(2ex,0ex) -- ++(1ex,-1ex)
            [rounded corners=2pt] -- ++(0ex,-3ex) -- cycle
            (element) ++(2ex,4ex) |- ++(1ex,-1ex);
        \draw (element) ++(3ex,0ex) node [fill=black, inner sep=0.5ex] {\texttt{AUX}};
    \end{scope}

    \draw [->, very thick] (compiler2) -- (log2);
    \draw [->, very thick] (compiler2) -- (aux2);
\end{tikzpicture}
\end{center}
\end{frame}





\subsection{El problema de {\lmr\BibTeX}}

\begin{frame}
\frametitle{Problemas de {\lmr\BibTeX}}

\pause

\begin{itemize}[<+->]
    \item Falta de soporte a Unicode.
    \item Los estilos de citación son limitados.
    \item No personalización.
\end{itemize}
\end{frame}





\section{\texttt{biblatex} y \texttt{biber}}

\begin{frame}
\frametitle{\texttt{biblatex} y \texttt{biber}}

Debido a estas limitantes, se crearon muchos paquetes que buscan solucionar todos estos problemas.

\pause

Uno de los más utilizados es el paquete es \texttt{biblatex}. Este nos añade aún más tipos de elementos para la base de datos \texttt{.bib}
\end{frame}





\begin{frame}[fragile]
\frametitle{La base de datos}

\begin{columns}
    \begin{column}{0.5\textwidth}
        \begin{minted}[escapeinside=||]{latex}
\usepackage{biblatex}
\addbibresource{|\textit{documento}|}|\pause|

\cite{TeXbook}|\pause|

\nocite{TeXbook}
\nocite{*}|\pause|

\printbibliography
        \end{minted}
    \end{column}
    \pause
    {\color{gray} \vline}
    \begin{column}{0.5\linewidth}
        \lmr
        \begin{enumerate}
            \setbeamerfont{enumerate item}{family=\lmr}
            \setbeamercolor{enumerate item}{fg=white}
            \setbeamertemplate{enumerate item}{[\insertenumlabel]}
            \item Donald Knuth. \textit{The {\TeX} book}. Stanford University, 1 edition, 1986.
        \end{enumerate}
    \end{column}
\end{columns}
\end{frame}





\subsection{Compilación}

\begin{frame}
\frametitle{Compilación}

\begin{center}
\begin{tikzpicture}
    \begin{scope}[local bounding box=tex]
        \useasboundingbox (-2ex,-2.5ex) rectangle (2ex,2.5ex);
        \draw [very thick]
            [rounded corners=2pt] (-1.5ex,-2ex) coordinate (element) -- ++(0ex,4ex)
            [rounded corners=1pt] -- ++(2ex,0ex) -- ++(1ex,-1ex)
            [rounded corners=2pt] -- ++(0ex,-3ex) -- cycle
            (element) ++(2ex,4ex) |- ++(1ex,-1ex);
        \draw (element) ++(3ex,0ex) node [fill=black, inner sep=0.5ex] {\texttt{TEX}};
    \end{scope}

    \begin{scope}[local bounding box=compiler, shift={([xshift=2cm]tex)}]
        \alt<6->{\node [black!75] {\lmr\LaTeX};}{\node {\lmr\LaTeX};}
    \end{scope}

    \alt<6->{\draw [->, very thick, black!75] (tex) -- (compiler);}{\draw [->, very thick] (tex) -- (compiler);}

    \begin{scope}[local bounding box=pdf, shift={([xshift=2cm]compiler)}]
        \useasboundingbox (-2ex,-2.5ex) rectangle (2ex,2.5ex);
        \alt<6->{
            \draw [very thick, black!75]
                [rounded corners=2pt] (-1.5ex,-2ex) coordinate (element) -- ++(0ex,4ex)
                [rounded corners=1pt] -- ++(2ex,0ex) -- ++(1ex,-1ex)
                [rounded corners=2pt] -- ++(0ex,-3ex) -- cycle
                (element) ++(2ex,4ex) |- ++(1ex,-1ex);
        }{
            \draw [very thick]
                [rounded corners=2pt] (-1.5ex,-2ex) coordinate (element) -- ++(0ex,4ex)
                [rounded corners=1pt] -- ++(2ex,0ex) -- ++(1ex,-1ex)
                [rounded corners=2pt] -- ++(0ex,-3ex) -- cycle
                (element) ++(2ex,4ex) |- ++(1ex,-1ex);
        }
        \alt<6->{
            \draw [black!75] (element) ++(3ex,0ex) node [fill=black, inner sep=0.5ex] {\texttt{PDF}};
        }{
            \draw (element) ++(3ex,0ex) node [fill=black, inner sep=0.5ex] {\texttt{PDF}};
        }
    \end{scope}

    \alt<6->{\draw [->, very thick, black!75] (compiler) -- (pdf);}{\draw [->, very thick] (compiler) -- (pdf);}

    \begin{scope}[local bounding box=log, shift={([yshift=2cm]pdf)}]
        \useasboundingbox (-2ex,-2.5ex) rectangle (2ex,2.5ex);
        \alt<6->{
            \draw [very thick, black!75]
                [rounded corners=2pt] (-1.5ex,-2ex) coordinate (element) -- ++(0ex,4ex)
                [rounded corners=1pt] -- ++(2ex,0ex) -- ++(1ex,-1ex)
                [rounded corners=2pt] -- ++(0ex,-3ex) -- cycle
                (element) ++(2ex,4ex) |- ++(1ex,-1ex);
        }{
            \draw [very thick]
                [rounded corners=2pt] (-1.5ex,-2ex) coordinate (element) -- ++(0ex,4ex)
                [rounded corners=1pt] -- ++(2ex,0ex) -- ++(1ex,-1ex)
                [rounded corners=2pt] -- ++(0ex,-3ex) -- cycle
                (element) ++(2ex,4ex) |- ++(1ex,-1ex);
        }
        \alt<6->{
            \draw [black!75] (element) ++(3ex,0ex) node [fill=black, inner sep=0.5ex] {\texttt{LOG}};
        }{
            \draw (element) ++(3ex,0ex) node [fill=black, inner sep=0.5ex] {\texttt{LOG}};
        }
    \end{scope}

    \begin{scope}[local bounding box=aux, shift={([yshift=-2cm]pdf)}]
        \useasboundingbox (-2ex,-2.5ex) rectangle (2ex,2.5ex);
        \alt<6->{
            \draw [very thick, black!75]
                [rounded corners=2pt] (-1.5ex,-2ex) coordinate (element) -- ++(0ex,4ex)
                [rounded corners=1pt] -- ++(2ex,0ex) -- ++(1ex,-1ex)
                [rounded corners=2pt] -- ++(0ex,-3ex) -- cycle
                (element) ++(2ex,4ex) |- ++(1ex,-1ex);
        }{
            \draw [very thick]
                [rounded corners=2pt] (-1.5ex,-2ex) coordinate (element) -- ++(0ex,4ex)
                [rounded corners=1pt] -- ++(2ex,0ex) -- ++(1ex,-1ex)
                [rounded corners=2pt] -- ++(0ex,-3ex) -- cycle
                (element) ++(2ex,4ex) |- ++(1ex,-1ex);
        }
        \alt<6->{
            \draw [black!75] (element) ++(3ex,0ex) node [fill=black, inner sep=0.5ex] {\texttt{AUX}};
        }{
            \draw (element) ++(3ex,0ex) node [fill=black, inner sep=0.5ex] {\texttt{AUX}};
        }
    \end{scope}

    \alt<6->{\draw [->, very thick, black!75] (compiler) -- (log);}{\draw [->, very thick] (compiler) -- (log);}
    \alt<6->{\draw [->, very thick, black!75] (compiler) -- (aux);}{\draw [->, very thick] (compiler) -- (aux);}

    \pause

    \begin{scope}[local bounding box=bib,shift={([yshift=2cm]tex)}]
        \useasboundingbox (-2ex,-2.5ex) rectangle (2ex,2.5ex);
        \draw [very thick]
            [rounded corners=2pt] (-1.5ex,-2ex) coordinate (element) -- ++(0ex,4ex)
            [rounded corners=1pt] -- ++(2ex,0ex) -- ++(1ex,-1ex)
            [rounded corners=2pt] -- ++(0ex,-3ex) -- cycle
            (element) ++(2ex,4ex) |- ++(1ex,-1ex);
        \draw (element) ++(3ex,0ex) node [fill=black, inner sep=0.5ex] {\texttt{BIB}};
    \end{scope}

    \alt<6->{\draw [->, very thick, shorten <=0.65em, black!75] (bib) -- (compiler);}{\draw [->, very thick, shorten <=0.65em] (bib) -- (compiler);}

    \pause

    \begin{scope}[local bounding box=biber, shift={([xshift=2cm]pdf)}]
        \alt<6->{\node [black!75] {\lmr Biber};}{\node {\lmr Biber};}
    \end{scope}

    \alt<6->{\draw [->, very thick, black!75] (aux) -- (biber);}{\draw [->, very thick] (aux) -- (biber);}

    \pause

    \begin{scope}[local bounding box=blg, shift={([xshift=2cm]biber)}]
        \useasboundingbox (-2ex,-2.5ex) rectangle (2ex,2.5ex);
        \alt<6->{
            \draw [very thick, black!75]
                [rounded corners=2pt] (-1.5ex,-2ex) coordinate (element) -- ++(0ex,4ex)
                [rounded corners=1pt] -- ++(2ex,0ex) -- ++(1ex,-1ex)
                [rounded corners=2pt] -- ++(0ex,-3ex) -- cycle
                (element) ++(2ex,4ex) |- ++(1ex,-1ex);
        }{
            \draw [very thick]
                [rounded corners=2pt] (-1.5ex,-2ex) coordinate (element) -- ++(0ex,4ex)
                [rounded corners=1pt] -- ++(2ex,0ex) -- ++(1ex,-1ex)
                [rounded corners=2pt] -- ++(0ex,-3ex) -- cycle
                (element) ++(2ex,4ex) |- ++(1ex,-1ex);
        }
        \alt<6->{
            \draw [black!75] (element) ++(3ex,0ex) node [fill=black, inner sep=0.5ex] {\texttt{BLG}};
        }{
            \draw (element) ++(3ex,0ex) node [fill=black, inner sep=0.5ex] {\texttt{BLG}};
        }
    \end{scope}

    \begin{scope}[local bounding box=bbl, shift={([yshift=-2cm]blg)}]
        \useasboundingbox (-2ex,-2.5ex) rectangle (2ex,2.5ex);
        \draw [very thick]
            [rounded corners=2pt] (-1.5ex,-2ex) coordinate (element) -- ++(0ex,4ex)
            [rounded corners=1pt] -- ++(2ex,0ex) -- ++(1ex,-1ex)
            [rounded corners=2pt] -- ++(0ex,-3ex) -- cycle
            (element) ++(2ex,4ex) |- ++(1ex,-1ex);
        \draw (element) ++(3ex,0ex) node [fill=black, inner sep=0.5ex] {\texttt{BBL}};
    \end{scope}

    \alt<6->{\draw [->, very thick, black!75] (biber) -- (bbl);}{\draw [->, very thick] (biber) -- (bbl);}
    \alt<6->{\draw [->, very thick, black!75] (biber) -- (blg);}{\draw [->, very thick] (biber) -- (blg);}

    \pause

    \begin{scope}[local bounding box=compiler2, shift={([xshift=2cm]blg)}]
        \node {\lmr\LaTeX};
    \end{scope}

    \pause

    \draw [->, very thick] (tex) -- (compiler2);
    \draw [->, very thick, rounded corners=5pt] (bib) -| (compiler2);
    \draw [->, very thick] (bbl) -- (compiler2);

    \pause

    \begin{scope}[local bounding box=pdf2, shift={([xshift=2cm]compiler2)}]
        \useasboundingbox (-2ex,-2.5ex) rectangle (2ex,2.5ex);
        \draw [very thick]
            [rounded corners=2pt] (-1.5ex,-2ex) coordinate (element) -- ++(0ex,4ex)
            [rounded corners=1pt] -- ++(2ex,0ex) -- ++(1ex,-1ex)
            [rounded corners=2pt] -- ++(0ex,-3ex) -- cycle
            (element) ++(2ex,4ex) |- ++(1ex,-1ex);
        \draw (element) ++(3ex,0ex) node [fill=black, inner sep=0.5ex] {\texttt{PDF}};
    \end{scope}

    \draw [->, very thick] (compiler2) -- (pdf2);

    \begin{scope}[local bounding box=log2, shift={([yshift=2cm]pdf2)}]
        \useasboundingbox (-2ex,-2.5ex) rectangle (2ex,2.5ex);
        \draw [very thick]
            [rounded corners=2pt] (-1.5ex,-2ex) coordinate (element) -- ++(0ex,4ex)
            [rounded corners=1pt] -- ++(2ex,0ex) -- ++(1ex,-1ex)
            [rounded corners=2pt] -- ++(0ex,-3ex) -- cycle
            (element) ++(2ex,4ex) |- ++(1ex,-1ex);
        \draw (element) ++(3ex,0ex) node [fill=black, inner sep=0.5ex] {\texttt{LOG}};
    \end{scope}

    \begin{scope}[local bounding box=aux2, shift={([yshift=-2cm]pdf2)}]
        \useasboundingbox (-2ex,-2.5ex) rectangle (2ex,2.5ex);
        \draw [very thick]
            [rounded corners=2pt] (-1.5ex,-2ex) coordinate (element) -- ++(0ex,4ex)
            [rounded corners=1pt] -- ++(2ex,0ex) -- ++(1ex,-1ex)
            [rounded corners=2pt] -- ++(0ex,-3ex) -- cycle
            (element) ++(2ex,4ex) |- ++(1ex,-1ex);
        \draw (element) ++(3ex,0ex) node [fill=black, inner sep=0.5ex] {\texttt{AUX}};
    \end{scope}

    \draw [->, very thick] (compiler2) -- (log2);
    \draw [->, very thick] (compiler2) -- (aux2);
\end{tikzpicture}
\end{center}
\end{frame}





\section{La base de datos}

\begin{frame}
\frametitle{La base de datos}

La base de datos de bibliografía (\texttt{.bib}) son totalemente independientes del documento (\texttt{.tex}).

Por lo que estas pueden ser reutilizadas en múltiples documentos.

Para {\lmr\BibTeX} solamente serán incluidos aquellas entradas que sean citadas en el documento.
\end{frame}





\subsection{Entradas}

\begin{frame}[fragile]
\frametitle{La base de datos}

Todas las entradas de la base de datos tiene la forma

\begin{minted}{bibtex}
@tipo{clave,
    parametro = {valor},
    parametro = "valor",
}
\end{minted}
\end{frame}





\begin{frame}[fragile]
\frametitle{La base de datos}

Se utilizará la siguiente nomenclatura

\begin{columns}
    \begin{column}{0.5\textwidth}
        \begin{minted}{bibtex}
@tipo{clave,
    param0 = {value},
    param1 = {value},
    ALTparam2 = {value},
    ALTparam3 = {value},
    OPTparam4 = {value},
    OPTparam5 = {value},
    OPTparam6 = {value},
    OPTparam7 = {value},
}
        \end{minted}
    \end{column}
    \pause
    \begin{column}{0.5\linewidth}
        \begin{minted}{bibtex}
@tipo{clave,
    param0 = {value},
    param1 = {value},
    param3 = {value},
    param4 = {value},
    param7 = {value},
}
        \end{minted}
    \end{column}
\end{columns}
\end{frame}





\begin{frame}[fragile]
\frametitle{La base de datos}
\framesubtitle{Artículo}

\begin{multicols}{2}
    \small
    \begin{minted}{bibtex}
@Article{key,
    author = {},
    title = {},
    journaltitle = {},
    date = {},
    OPTkey = {},
    OPTsubtitle = {},
    OPTeditor = {},
    OPTissuetitle = {},
    OPTlanguage = {},
    OPToriglanguage = {},
    OPTseries = {},
    OPTvolume = {},
    OPTnumber = {},
    OPTeid = {},
    OPTissue = {},
    OPTpages = {},
    OPTversion = {},
    OPTnote = {},
    OPTissn = {},
    OPTaddendum = {},
    OPTpubstate = {},
    OPTdoi = {},
    OPTurl = {},
    OPTurldate = {},
    OPTannote = {},
}
    \end{minted}
\end{multicols}
\end{frame}





\begin{frame}[fragile]
\frametitle{La base de datos}
\framesubtitle{Libro}

\begin{multicols}{2}
    \small
    \begin{minted}{bibtex}
@Book{key,
    author = {},
    title = {},
    date = {},
    OPTkey = {},
    OPTeditor = {},
    OPTsubtitle = {},
    OPTlanguage = {},
    OPToriglanguage = {},
    OPTvolume = {},
    OPTpart = {},
    OPTedition = {},
    OPTvolumes = {},
    OPTseries = {},
    OPTnumber = {},
    OPTpublisher = {},
    OPTlocation = {},
    OPTisbn = {},
    OPTeid = {},
    OPTchapter = {},
    OPTpages = {},
    OPTpagetotal = {},
    OPTdoi = {},
    OPTurl = {},
    OPTurldate = {},
    OPTannote = {},
}
    \end{minted}
\end{multicols}
\end{frame}





\begin{frame}[fragile]
\frametitle{La base de datos}
\framesubtitle{Cuadernillo, Folleto o Panfleto}

\begin{multicols}{2}
    \small
    \begin{minted}{bibtex}
@Booklet{key,
    ALTauthor = {},
    ALTeditor = {},
    title = {},
    date = {},
    OPTkey = {},
    OPTsubtitle = {},
    OPTlanguage = {},
    OPThowpublished = {},
    OPTtype = {},
    OPTnote = {},
    OPTlocation = {},
    OPTeid = {},
    OPTchapter = {},
    OPTpages = {},
    OPTpagetotal = {},
    OPTaddendum = {},
    OPTpubstate = {},
    OPTdoi = {},
    OPTurl = {},
    OPTurldate = {},
    OPTannote = {},
}
    \end{minted}
\end{multicols}
\end{frame}






\begin{frame}[fragile]
\frametitle{La base de datos}
\framesubtitle{Conferencia}

\begin{multicols}{2}
    \small
    \begin{minted}{bibtex}
@Conference{key,
    author = {},
    title = {},
    booktitle = {},
    year = {},
    OPTkey = {},
    OPTeditor = {},
    OPTvolume = {},
    OPTnumber = {},
    OPTseries = {},
    OPTpages = {},
    OPTaddress = {},
    OPTorganization = {},
    OPTpublisher = {},
    OPTnote = {},
    OPTannote = {},
}
    \end{minted}
\end{multicols}
\end{frame}






\begin{frame}[fragile]
\frametitle{La base de datos}
\framesubtitle{Dentro de un Libro}

\begin{multicols}{2}
    \small
    \begin{minted}{bibtex}
@InBook{key,
    title = {},
    date = {},
    author = {},
    booktitle = {},
    OPTkey = {},
    OPTbookauthor = {},
    OPTeditor = {},
    OPTsubtitle = {},
    OPTlanguage = {},
    OPToriglanguage = {},
    OPTvolume = {},
    OPTpart = {},
    OPTedition = {},
    OPTvolumes = {},
    OPTseries = {},
    OPTnumber = {},
    OPTnote = {},
    OPTpublisher = {},
    OPTlocation = {},
    OPTisbn = {},
    OPTeid = {},
    OPTchapter = {},
    OPTpages = {},
    OPTdoi = {},
    OPTurl = {},
    OPTurldate = {},
    OPTannote = {},
}
    \end{minted}
\end{multicols}
\end{frame}





\begin{frame}[fragile]
\frametitle{La base de datos}
\framesubtitle{En una colección}

\begin{multicols}{2}
    \small
    \begin{minted}{bibtex}
@InCollection{key,
    author = {},
    title = {},
    date = {},
    OPTcrossref = {},
    OPTkey = {},
    OPTbooktitle = {},
    OPTeditor = {},
    OPTsubtitle = {},
    OPTlanguage = {},
    OPToriglanguage = {},
    OPTvolume = {},
    OPTpart = {},
    OPTedition = {},
    OPTvolumes = {},
    OPTseries = {},
    OPTnumber = {},
    OPTnote = {},
    OPTpublisher = {},
    OPTlocation = {},
    OPTisbn = {},
    OPTeid = {},
    OPTchapter = {},
    OPTpages = {},
    OPTdoi = {},
    OPTurl = {},
    OPTurldate = {},
    OPTannote = {},
}
    \end{minted}
\end{multicols}
\end{frame}





\begin{frame}[fragile]
\frametitle{La base de datos}
\framesubtitle{En un artículo en una conferencia}

\begin{multicols}{2}
    \small
    \begin{minted}{bibtex}
@InProceedings{key,
    author = {},
    title = {},
    date = {},
    OPTcrossref = {},
    OPTkey = {},
    OPTbooktitle = {},
    OPTeditor = {},
    OPTsubtitle = {},
    OPTeventtitle = {},
    OPTeventdate = {},
    OPTvenue = {},
    OPTlanguage = {},
    OPTvolume = {},
    OPTpart = {},
    OPTvolumes = {},
    OPTseries = {},
    OPTnumber = {},
    OPTnote = {},
    OPTorganization = {},
    OPTpublisher = {},
    OPTlocation = {},
    OPTchapter = {},
    OPTpages = {},
    OPTdoi = {},
    OPTurl = {},
    OPTurldate = {},
    OPTannote = {},
}
    \end{minted}
\end{multicols}
\end{frame}






\begin{frame}[fragile]
\frametitle{La base de datos}
\framesubtitle{Manuales o Documentación Técnica}

\begin{multicols}{2}
    \small
    \begin{minted}{bibtex}
@Manual{key,
    ALTauthor = {},
    ALTeditor = {},
    title = {},
    date = {},
    OPTkey = {},
    OPTsubtitle = {},
    OPTlanguage = {},
    OPTedition = {},
    OPTtype = {},
    OPTseries = {},
    OPTnumber = {},
    OPTversion = {},
    OPTnote = {},
    OPTorganization = {},
    OPTpublisher = {},
    OPTlocation ={},
    OPTisbn = {},
    OPTeid = {},
    OPTchapter = {},
    OPTpages = {},
    OPTpagetotal = {},
    OPTaddendum = {},
    OPTpubstate = {},
    OPTdoi = {},
    OPTurl = {},
    OPTurldate = {},
    OPTannote = {},
}
    \end{minted}
\end{multicols}
\end{frame}





\begin{frame}[fragile]
\frametitle{La base de datos}
\framesubtitle{Tesis de Maestría}

\begin{multicols}{2}
    \small
    \begin{minted}{bibtex}
@MastersThesis{key,
    author = {},
    title = {},
    school = {},
    year = {},
    OPTkey = {},
    OPTtype = {},
    OPTaddress = {},
    OPTmonth = {},
    OPTnote = {},
    OPTannote = {},
}
    \end{minted}
\end{multicols}
\end{frame}





\begin{frame}[fragile]
\frametitle{La base de datos}
\framesubtitle{Tesis de Doctorado}

\begin{multicols}{2}
    \small
    \begin{minted}{bibtex}
@PhdThesis{key,
    author = {},
    title = {},
    school = {},
    year = {},
    OPTkey = {},
    OPTtype = {},
    OPTaddress = {},
    OPTmonth = {},
    OPTnote = {},
    OPTannote = {}
}
    \end{minted}
\end{multicols}
\end{frame}





\begin{frame}[fragile]
\frametitle{La base de datos}
\framesubtitle{Misceláneo}

\begin{multicols}{2}
    \small
    \begin{minted}{bibtex}
@Misc{key,
    ALTauthor = {},
    ALTeditor = {},
    title = {},
    date = {},
    OPTkey = {},
    OPTsubtitle = {},
    OPTtitleaddon = {},
    OPTlanguage =  {},
    OPThowpublished = {},
    OPTtype = {},
    OPTversion = {},
    OPTnote = {},
    OPTorganization = {},
    OPTlocation = {},
    OPTdoi = {},
    OPTurl = {},
    OPTurldate = {},
    OPTannote = {},
}
    \end{minted}
\end{multicols}
\end{frame}





\begin{frame}[fragile]
\frametitle{La base de datos}
\framesubtitle{Minuta o Acta de una conferencia}

\begin{multicols}{2}
    \small
    \begin{minted}{bibtex}
@Proceedings{key,
    title = {},
    date = {},
    OPTkey = {},
    OPTeditor = {},
    OPTsubtitle = {},
    OPTeventtitle = {},
    OPTeventdate = {},
    OPTvenue = {},
    OPTlanguage = {},
    OPTvolume = {},
    OPTpart = {},
    OPTvolumes = {},
    OPTseries = {},
    OPTnumber = {},
    OPTnote = {},
    OPTorganization = {},
    OPTpublisher = {},
    OPTlocation = {},
    OPTisbn = {},
    OPTeid = {},
    OPTchapter = {},
    OPTpages = {},
    OPTpagetotal = {},
    OPTdoi = {},
    OPTurl = {},
    OPTurldate = {},
    OPTannote = {},
}
    \end{minted}
\end{multicols}
\end{frame}





\begin{frame}[fragile]
\frametitle{La base de datos}
\framesubtitle{Reporte Técnico}

\begin{multicols}{2}
    \small
    \begin{minted}{bibtex}
@TechReport{key,
  author = {},
  title = {},
  institution = {},
  year = {},
  OPTkey = {},
  OPTtype = {},
  OPTnumber = {},
  OPTaddress = {},
  OPTmonth = {},
  OPTnote = {},
  OPTannote = {},
}
    \end{minted}
\end{multicols}
\end{frame}





\begin{frame}[fragile]
\frametitle{La base de datos}
\framesubtitle{Sin publicar}

\begin{multicols}{2}
    \small
    \begin{minted}{bibtex}
@Unpublished{key,
    author = {},
    title = {},
    date = {},
    OPTkey = {},
    OPTsubtitle = {},
    OPTtype = {},
    OPTeventtitle = {},
    OPTeventtitleaddon = {},
    OPTeventdate = {},
    OPTvenue = {},
    OPTlanguage = {},
    OPThowpublished = {},
    OPTnote = {},
    OPTlocation = {},
    OPTisbn = {},
    OPTdoi = {},
    OPTurl = {},
    OPTurldate = {},
    OPTannote = {},
}
    \end{minted}
\end{multicols}
\end{frame}




\begin{frame}[allowframebreaks]
\frametitle{Referencias}
\nocite{*}
\printbibliography
\end{frame}
\end{document}
